\section{Number, Math, dan Pembulatan}

Tipe \code{Number} di JavaScript menggunakan IEEE 754 double-precision float (64-bit). Ini berarti semua angka (integer dan desimal) disimpan dalam format yang sama. Konsekuensinya: \code{0.1 + 0.2 !== 0.3} (masalah floating-point). Untuk perhitungan uang, gunakan integer (sen) atau library seperti Decimal.js. \code{Number.isInteger()}, \code{Number.isFinite()}, dan \code{Number.isNaN()} memeriksa sifat angka \cite{jsInfo}.

\begin{table}[htbp]
\centering
\begin{tabular}{|P{3.5cm}|P{4cm}|P{4.5cm}|}
\hline
\textbf{Metode/Properti} & \textbf{Fungsi} & \textbf{Contoh} \\
\hline
\code{Math.round(x)} & Bulatkan terdekat & \code{Math.round(4.5)} $\rightarrow$ 5 \\
\hline
\code{Math.floor(x)} & Bulatkan ke bawah & \code{Math.floor(4.9)} $\rightarrow$ 4 \\
\hline
\code{Math.ceil(x)} & Bulatkan ke atas & \code{Math.ceil(4.1)} $\rightarrow$ 5 \\
\hline
\code{Math.random()} & Acak 0--1 & \code{Math.random()} \\
\hline
\code{Math.max(...vals)} & Nilai terbesar & \code{Math.max(1,5,3)} $\rightarrow$ 5 \\
\hline
\code{Number.parseInt(s)} & Parse ke integer & \code{parseInt("42px")} $\rightarrow$ 42 \\
\hline
\code{.toFixed(n)} & Format desimal & \code{(3.14159).toFixed(2)} $\rightarrow$ \code{"3.14"} \\
\hline
\end{tabular}
\caption{Metode Math dan Number yang sering digunakan}
\end{table}

\begin{konsep}
\textbf{Random Integer dalam Rentang.} Untuk menghasilkan integer acak antara \code{min} dan \code{max} (inklusif): \code{Math.floor(Math.random() * (max - min + 1)) + min}. \code{Math.random()} menghasilkan float 0 $\leq$ x $<$ 1; kalikan dengan rentang, bulatkan ke bawah, lalu tambahkan minimum \cite{mdnJsGuide}.
\end{konsep}

